\section{讨论与限制}

\subsection{“突围/收编”的可测边界}

结果支持一种比标题更克制的解释：思想相关人物在关系图中确有较高比例的师承边，但他们的大多数显式关系仍通向朝廷权力。所谓“突围”可由师生同门占比和1508--1528年的传播事件簇近似；所谓“收编”只能由权力关系占比、1579年关闭书院和1593年讲学入党争近似。两者都是分析者施加的框架，不是数据直接吐出的标签。本文不声称测得思想自由度，也不把书院关闭与所有思想控制画等号。

\subsection{首要效度威胁：没有思想文本本身}

本文最大的限制不是样本量，而是材料性质。语料只有通俗史著中的叙事提及，没有《传习录》原文、门人问答、书院规约、奏疏或讲会记录。于是“知行合一”12次只表示这个字符串在叙事中出现12次，不能说明命题被讨论了12次，更不能衡量读者接受。叙事作者把王守仁写成能平叛的行动者，也会同时提高其政治事件数与关系度数；这种写法本身就是研究对象，却使思想内容与叙事情节无法分离。

\subsection{关系抽取与分组内造}

关系类别由项目抽取体系预先定义。“师生同门”把严格师承、同门、门生与思想认同放在一类；四类“权力关系”则同时包含制度隶属、攻讦、冲突与联盟。它们适合做结构描述，不适合推断情感或动机。思想组名单也由同一语料的角色、事件与关系标记确定，存在同源自证。剔除58条思想组推导边后，师承占比从18.2\%升到30.0\%，说明自动共现对构成有显著影响；但稳健性方向未翻转。

\subsection{外推与复现}

本研究只是一部书、一个抽取版本的案例。要讨论明代思想传播本身，下一步至少应加入思想原典与不同史著，区分著述、讲会、科举、官僚任迁四种关系，并让不同标注者独立复核类别。本文公开统计底稿、筛选规则与两份作图脚本；版权文本不公开，但每处书内证据都保留逐字命中次数。引用只使用核验通过的条目，语料书籍按无 DOI 版本信息列示。
